\section{Preguntas sin resolver, a fondo: las cinco rutas con más probabilidad de divergir antes de 2027}

\subsection{Ruta uno: ¿se convertirá el inference-time compute en la nueva normalidad?}

La postura de la entrevista sobre el inference-time scaling es positiva pero cautelosa. Es positiva porque puede elevar rápidamente la calidad en tareas de alta dificultad; es cautelosa porque de manera inherente eleva el costo y hace la experiencia más variable. La pregunta más importante antes de 2027 no es «si se puede hacer», sino «en qué tareas vale la pena hacerlo».

\begin{importantbox}{Un criterio de decisión aplicable}
Si una tarea cumple las dos condiciones siguientes, se debe priorizar activar el modo de presupuesto de inferencia alto:
\begin{itemize}[leftmargin=1.5em]
    \item el costo de un error es mayor que el costo de la latencia (por ejemplo, código de producción, documentos de cumplimiento normativo);
    \item existe un mecanismo de validación claro que permita determinar si «calcular más produce un beneficio cuantificable».
\end{itemize}
En caso contrario, se sigue por defecto la ruta rápida y se eleva el nivel según sea necesario mediante un enrutador.
\end{importantbox}

\subsection{Ruta dos: la división de trabajo a largo plazo entre Open-weight y Closed API}

En la entrevista, las dos rutas coexisten y se refuerzan mutuamente. Open-weight impulsa la transparencia de la investigación y la flexibilidad de despliegue; Closed API impulsa la integración profunda de herramientas y la unificación del producto. Lo más probable a futuro es una coexistencia por capas, no el triunfo de una sola parte.

\begin{warningbox}{El error más común en las decisiones organizacionales}
Tratar la «apertura» como un problema puramente ideológico y no como un problema de negocio. Para una empresa, primero debe preguntarse: qué tan altos son, respectivamente, la sensibilidad de los datos, las restricciones de despliegue, la necesidad de interpretabilidad y el riesgo de dependencia de un proveedor, y solo después decidir la proporción de adopción.
\end{warningbox}

\subsection{Ruta tres: si la capacidad del post-entrenamiento puede transferirse a tareas no verificables}

La fortaleza de RLVR proviene de ser «verificable». La pregunta realmente pendiente es si, después de reforzarse en una gran cantidad de tareas verificables, el modelo obtendrá beneficios estables en tareas de verificación débil como la creatividad, la estrategia y las hipótesis científicas. La entrevista no ofrece una conclusión, pero sí una dirección de investigación: recompensas de proceso más sólidas y entrenamiento mixto multitarea.

\begin{knowledgebox}{Sugerencia de investigación aplicable}
Se puede usar «tareas semiverificables» como puente: primero se definen subobjetivos locales verificables, y después se evalúa la mejora de la tarea en su conjunto. Esto es más experimentable y más reproducible que preguntar directamente si RLVR puede mejorar la creatividad.
\end{knowledgebox}

\subsection{Ruta cuatro: ¿el human-in-the-loop se reducirá o se ampliará?}

A corto plazo, la mejora de eficiencia que aporta la IA reducirá la proporción de participación humana; a mediano y largo plazo, la participación humana en los eslabones de mayor valor se concentrará más y será más costosa. La discusión de la entrevista sobre «voice», «agency» y «meaning» apunta, en esencia, a este mismo punto: el papel humano se reduce en número mientras aumenta en densidad.

\begin{importantbox}{El mínimo indispensable en el diseño de la colaboración entre humano y máquina}
El sistema debe responder tres preguntas de responsabilidad:
\begin{itemize}[leftmargin=1.5em]
    \item quién define el objetivo y asume las consecuencias;
    \item quién revisa los resultados intermedios clave;
    \item quién tiene la autoridad de frenar de emergencia en caso de falla.
\end{itemize}
Si estos tres puntos no son claros, por más alta que sea la tasa de automatización, se convertirá en un riesgo operativo.
\end{importantbox}

\subsection{Ruta cinco: si la velocidad de la gobernanza puede mantenerse a la par de la velocidad tecnológica}

Desde las licencias de datos hasta las políticas de código abierto y la seguridad psicológica, los temas de gobernanza ya atraviesan todo el proceso completo. El juicio realista que ofrece la entrevista es este: la gobernanza no se pondrá al día de manera natural; es necesario integrar de forma sistemática los requisitos de gobernanza en el proceso de desarrollo, en lugar de aplicarlos como un parche posterior.

\begin{warningbox}{En el contexto de la IA, «lanzar primero y corregir después» cuesta más caro}
Los sistemas de modelos tienen una naturaleza de propagación rápida, y las conductas erróneas se replican a gran escala. Si el proceso de gobernanza llega después del ritmo de lanzamiento, el costo de corrección aumenta de forma exponencial, e incluye riesgo de marca, riesgo legal y riesgo para la confianza social.
\end{warningbox}

\subsection{Resumen de la sección}

Antes de 2027, lo más importante a seguir no es un único modelo ganador, sino cuál de estas cinco rutas es la primera en formar un ciclo estable: presupuesto de inferencia, división de la apertura, transferencia del aprendizaje, responsabilidad entre humano y máquina, e integración de la gobernanza. Solo donde se forme ese ciclo aparecerá una ventaja competitiva verdaderamente sostenida.

